% DOCUMENT CLASS %
\documentclass[a4paper, 12pt]{extarticle}   % for A4 Paper, 12pt fontsize and article spacing
\usepackage[left=1.5cm,right=1.5cm,top=2cm,bottom=2cm]{geometry}    % Make the page a little bigger

% PACKAGES%
\usepackage{amsfonts}   % For a basic mathfont (many will be replaced by stix)
\usepackage{amsmath}    % For basic math symbols
\usepackage{bbm}        % For \mathbbm (better version of \mathbb)
\usepackage{mathrsfs}   % For \mathscr
\usepackage{hyperref}   % For footnotes
\usepackage{csquotes}   % For \textquote{}
\usepackage{IEEEtrantools}  % For better alignments
\usepackage{tikz}   % For a macro
\usetikzlibrary{arrows.meta, positioning, calc}	% For protocol diagrams
\usepackage{listings} % For code
\usepackage{xcolor}   % For ... code
\usepackage{microtype}  % For better paragraph formatting
\usepackage{amsthm}   % For custom environments

% FONT CONFIGURATION %
\usepackage[T1]{fontenc}
\usepackage{palatino}

% BASIC SETS %
\newcommand*{\R}{\mathbbm R}    % Set of real numbers
\newcommand*{\N}{\mathbbm N}    % Set of natural numbers (beginning at 0)
\newcommand*{\Z}{\mathbbm Z}    % Set of integers
\newcommand*{\Q}{\mathbbm Q}    % Set of rational numbers

% MACROS %
\newcommand*{\qedbox}{\null\nobreak\hfill\ensuremath{\square}} % Macro for a QED-Box
% for a funny ¯\_(ツ)_/¯-Emoji
\newcommand{\shrug}[1][]{%
    \begin{tikzpicture}[baseline,x=0.8\ht\strutbox,y=0.8\ht\strutbox,line width=0.125ex,#1]
    \def\arm{(-2.5,0.95) to (-2,0.95) (-1.9,1) to (-1.5,0) (-1.35,0) to (-0.8,0)};
    \draw \arm;
    \draw[xscale=-1] \arm;
    \def\headpart{(0.6,0) arc[start angle=-40, end angle=40,x radius=0.6,y radius=0.8]};
    \draw \headpart;
    \draw[xscale=-1] \headpart;
    \def\eye{(-0.075,0.15) .. controls (0.02,0) .. (0.075,-0.15)};
    \draw[shift={(-0.3,0.8)}] \eye;
    \draw[shift={(0,0.85)}] \eye;
    % draw mouth
    \draw (-0.1,0.2) to [out=15,in=-100] (0.4,0.95); 
    \end{tikzpicture}
}

% BASIC CONFIGURATION %
\pagestyle{plain}
\allowdisplaybreaks

% CODE CONFIGURATIONS %
\definecolor{codepurple}{HTML}{7f0055}
\definecolor{comment}{RGB}{0,128,0}
\definecolor{number}{RGB}{9,136,90}
\definecolor{string}{RGB}{163,21,21}

\lstset{
    language=python,
    emph={*, do, od, for, if, fi, then, else, to, def, in, forall, exists},
    firstnumber=1,
    basicstyle=\small\ttfamily,
    breaklines=true,
    breakatwhitespace=false,
    frame=tb,
    captionpos=b,
    %float=t,
    %floatplacement=t,
    abovecaptionskip=0.5cm,
    numbers=left,
    numberstyle=\tiny,
    keywordstyle=\bfseries\color{codepurple},
    emphstyle=\bfseries\color{codepurple},
    mathescape=true,
    keepspaces=true,
    showspaces=false,
    showstringspaces=false,
    showtabs=false,
    stringstyle=\color{string},
    commentstyle=\itshape\color{comment},
    upquote=true,
    texcl=true
}

% EXERCISE ENVIRONMENT %
\theoremstyle{definition}
\newtheorem{exercise}{Exercise}[section]
\setcounter{section}{1} % Don't use \section and this works naturally

% TITLE CONFIGURATION %
\newcommand{\ex}{0}

\title{Data Visualization \\ \small Exercise Sheet \ex}
\date{Winter Term 2025/26}
\author{
    Henri Heyden \\%
    \small stu240825%
    \and
    Julia Schröder \\%
    \small stu252018%
    \and
    Rebecca Ziebell \\%
    \small stu250567%
}

\renewcommand{\ex}{7}
\setcounter{section}{\ex}

\begin{document}
    \maketitle
    
    \begin{exercise}[Warmup] \hfill
        \begin{enumerate}
            \item[a)] The steps to-do are as follows:
            \begin{enumerate}
                \item[1.] Centre/normalize\footnote{Standard normalization is generally not needed, PCA only requires centred points.} points in matrix \(X \in \R^{d \times n}\) if the mean dimensional components are not zero.
                \item[2.] Calculate the covariance matrix \(C_X := 1/n \cdot X^TX\).\footnote{The scalar is not useful, we only need a matrix, which is proportional to the \emph{actual} covariance matrix.}
                \item[3.] Find the eigenvalues of \(C_X\) by solving for the roots of the characteristic polynomial of \(C\).
                Do this by solving for all \(\lambda\) in \(0 = \det(C_X - \lambda I_d)\), where \(I_d\) is the identity matrix of dimension \(d\).
                \item[4.] Find the eigenvectors of \(C_X\) by finding \(v \in \R^d\), such that \((0)^d = (C_X - \lambda I_d) \times v\) for all eigenvalues \(\lambda\).
                \item[5.] All orthogonal directions in \(X\) in decreasing variance are found by ordering the eigenvectors by decreasing absolute corresponding eigenvalue.
                \item[6.] The centred dataset can be projected on a dimensionality reduced space by picking the most variant directions, normalizing them, and putting them column-wise a matrix \(V\), and multiplying \(P = X \times P\).
            \end{enumerate}

            \item[b)]
            We get the centred matrix
            
            \[\begin{bmatrix}
                -1 & -1 \\
                1 & 1 \\
                -1 & 1 \\
                1 & -1
            \end{bmatrix}.\]

            Thus, we get the covariance matrix

            \[\begin{bmatrix}
                1 & 0 \\
                0 & 1
            \end{bmatrix}.\]

            \item[c)]
            We firstly calculate the eigenvalues

            \[0 = \det
            \begin{bmatrix}
                -5 - \lambda && 2 \\
                -7 && 4 - \lambda
            \end{bmatrix}
            = x^2 + x - 6.\]

            This polynomial has solutions \(\lambda_1 = -3, \lambda_2 = 2\).
            For the eigenvector of \(\lambda_1\) we have to find non-trivial \(v_1\), such that
            
            \[
            \begin{bmatrix}
                -2 & 2 \\
                -7 & 7
            \end{bmatrix}
            \times
            \begin{bmatrix}
                v_1^1 \\
                v_1^2
            \end{bmatrix}
            =
            \begin{bmatrix}
            0 \\
            0
            \end{bmatrix}.
            \]

            We can pick \(v_1 = \begin{bmatrix} 1 & 1 \end{bmatrix}^T\) which is our first eigenvector.
            For the second eigenvalue, we have to solve 

            \[
            \begin{bmatrix}
                -7 & 2 \\
                -7 & 2
            \end{bmatrix}
            \times
            \begin{bmatrix}
                v_2^1 \\
                v_2^2
            \end{bmatrix}
            =
            \begin{bmatrix}
            0 \\
            0
            \end{bmatrix}.
            \]

            A solution, the second eigenvector, is found using \(v_2 = \begin{bmatrix} 2 & 7 \end{bmatrix}^T\).

            \item[d)]
            The matrix has the characteristic polynomial of \(-x^3 + 25x^2 - 200x + 500\) which has (real) roots \(\lambda_1 = 10, \lambda_2 = 5\). Thus for the first eigenvector we have the following linear equation system.

            \begin{equation*}
                \begin{IEEEeqnarraybox}[\IEEEeqnarraystrutmode]{l'v;r'r'r;v'r}
                    \text{I}    && -5v_1^1 & -10v_2^1 & -5v_3^1 && 0 \\
                    \text{II}   &&  2v_1^1 & -4v_2^1  &  2v_3^1 && 0 \\
                    \text{III}  && -4v_1^1 & -8v_2^1  & -4v_3^1 && 0
                \end{IEEEeqnarraybox}
            \end{equation*}

            From \(\text{I} + 5/4 \cdot \text{III} = 0\) we get \(\exists a \in \R: av_1^1 + av_3^1 = 0\).
            Therefore, we know \(v_1^1 = -v_3^1\).
            From this we can transform II into \(0v_1^1 - 4v_2^1 = 0\) and eliminate \(v_2^1\).
            Thus, we obtain \(v_1 = \begin{bmatrix}1 & 0 & -1\end{bmatrix}^T\).

            For the second eigenvalue we obtain one eigenvector \(v_2 = \begin{bmatrix}-5 & 2 & -4\end{bmatrix}^T\) by substitution.

            \item[e)]
            Using the previously described steps, we obtain a factor of the covariance matrix:

            \[
            \begin{bmatrix}
                21 & 7 & -24 \\
                7 & 532/3 & 17 \\
                -24 & 17 & 31
            \end{bmatrix}.
            \]

            We obtain the largest eigenvalue \(\lambda_1 \approx 178\), with eigenvector \(v_1 = \begin{bmatrix}0.028 & 0.994 & 0.11\end{bmatrix}^T\).

            To project the points \(p_1\) and \(p_2\) to the direction, we have to center them again first to obtain \(\hat p_1 = \begin{bmatrix}1 & 13/3 & 3\end{bmatrix}\) and \(\hat p_2 = \begin{bmatrix}4 & 130/3 & 22\end{bmatrix}\).

            Finally, we obtain the first projection \(4.69 \approx \hat p_1 v_1\) and second projection \(45.89 \approx v_1 \hat p_2\).
        \end{enumerate}
    \end{exercise}

    \begin{exercise}[Data]
    \end{exercise}

    \begin{exercise}{PCA}
    \end{exercise}
    
    \begin{exercise}[PCA -- SKLEARN]
    \end{exercise}
\end{document}
